\paragraph{window-6 三轴中心约束、Pati--Salam 连通实现与标准模型包络的中心 2-挠兼容性}
在 window-6 的 boundary 读出中，三轴独立 sheet-flip 给出一个抽象中心寄存器
\[
\mathcal C_{\mathrm{bdry}}\cong (\ZZ/2\ZZ)^3.
\]
本段讨论其与紧连通 Lie 包络在中心 2-挠层面的可实现条件。

\begin{theorem}[标准模型连通包络对三轴中心寄存器的不可实现性]\label{thm:xi-window6-sm-envelope-center-2torsion-obstruction}
设 \(G\) 为紧连通 Lie 群，满足
\[
\Lie(G)\cong \mathfrak u(1)\oplus\mathfrak{su}(2)\oplus\mathfrak{su}(3).
\]
则
\[
\dim_{\FF_2} Z(G)[2]\le 2.
\]
从而不存在单同态
\[
\mathcal C_{\mathrm{bdry}}\cong (\ZZ/2\ZZ)^3\hookrightarrow Z(G)[2].
\]
\end{theorem}
\begin{proof}
由紧连通群分类，可取
\[
G\cong \widetilde G/\Gamma,\qquad
\widetilde G:=U(1)\times SU(2)\times SU(3),
\]
其中 \(\Gamma\subset Z(\widetilde G)\) 为有限中心子群。记
\[
A:=Z(\widetilde G)\cong U(1)\times \ZZ_2\times \ZZ_3,
\qquad
Z(G)\cong A/\Gamma.
\]
设 \(\pi:A\to \ZZ_2\times \ZZ_3\) 为自然投影，则有短正合列
\[
1\to U(1)\to A\to \ZZ_2\times \ZZ_3\to 1.
\]
对 \(\Gamma\) 取商后得到
\[
1\to T\to A/\Gamma\to Q\to 1,
\]
其中 \(T\cong U(1)\)（一维紧环面），\(Q\) 为 \((\ZZ_2\times \ZZ_3)\) 的商群。

记 \(E:=A/\Gamma\)。考虑 \(E[2]\to Q[2]\)。对任意 \(q\in Q[2]\)，其纤维中两点之差落在 \(T[2]\)。
故每个纤维大小至多 \(|T[2]|=2\)，并且
\[
|Q[2]|\le |(\ZZ_2\times \ZZ_3)[2]|=2.
\]
从而
\[
|E[2]|\le 2\cdot 2=4,\qquad
\dim_{\FF_2}E[2]\le 2.
\]
即 \(\dim_{\FF_2} Z(G)[2]\le 2\)。因此 \((\ZZ/2\ZZ)^3\) 不能嵌入 \(Z(G)[2]\)。证毕。
\end{proof}

\begin{theorem}[三轴 canonical sheet-flip 保真下降下的 Pati--Salam 连通实现唯一性]\label{thm:xi-window6-ps-faithful-canonical-descent-unique}
设
\[
\widetilde G_{\mathrm{PS}}:=SU(4)\times SU(2)_L\times SU(2)_R,\qquad
G:=\widetilde G_{\mathrm{PS}}/H,
\]
其中 \(H\subset Z(\widetilde G_{\mathrm{PS}})\) 为有限中心子群。记 canonical involution 子群
\[
\mathcal C_{\mathrm{can}}
:=Z(SU(4))[2]\times Z(SU(2)_L)\times Z(SU(2)_R)
\cong (\ZZ/2\ZZ)^3.
\]
若商映射 \(q:\widetilde G_{\mathrm{PS}}\to G\) 在 \(\mathcal C_{\mathrm{can}}\) 上限制为单同态
\[
q|_{\mathcal C_{\mathrm{can}}}:\mathcal C_{\mathrm{can}}\hookrightarrow Z(G)[2],
\]
则必有 \(H=\{e\}\)，因而
\[
G\cong SU(4)\times SU(2)_L\times SU(2)_R.
\]
\end{theorem}
\begin{proof}
有
\[
Z(\widetilde G_{\mathrm{PS}})
\cong \ZZ_4\times \ZZ_2\times \ZZ_2.
\]
取任意非平凡 \(H\subset Z(\widetilde G_{\mathrm{PS}})\)。任取 \(h\in H\setminus\{e\}\)。
若 \(\ord(h)=2\)，则 \(h\in \mathcal C_{\mathrm{can}}\)；
若 \(\ord(h)=4\)，则 \(h^2\neq e\) 且 \(h^2\in \mathcal C_{\mathrm{can}}\)。
故总有
\[
H\cap \mathcal C_{\mathrm{can}}\neq \{e\}.
\]
而 \(q\) 在 \(H\) 上恒等于单位元，故 \(q|_{\mathcal C_{\mathrm{can}}}\) 的核至少含
\(H\cap \mathcal C_{\mathrm{can}}\) 的非平凡元，不能单射，矛盾。
故 \(H\) 只能为平凡子群。证毕。
\end{proof}

\begin{theorem}[Pati--Salam 到标准模型包络的中心 2-挠压缩不可避免]\label{thm:xi-window6-ps-to-sm-center2-kernel}
设
\[
G_{\mathrm{PS}}:=SU(4)\times SU(2)_L\times SU(2)_R,
\]
且 \(G_{\mathrm{SM}}\) 为任意紧连通 Lie 群，满足
\(
\Lie(G_{\mathrm{SM}})\cong \mathfrak u(1)\oplus\mathfrak{su}(2)\oplus\mathfrak{su}(3)
\)。
若连续群同态
\(
f:G_{\mathrm{PS}}\to G_{\mathrm{SM}}
\)
满足中心兼容条件
\[
f\bigl(\mathcal C_{\mathrm{can}}\bigr)\subseteq Z(G_{\mathrm{SM}})[2],
\]
则限制映射
\[
f_*:\mathcal C_{\mathrm{can}}\longrightarrow Z(G_{\mathrm{SM}})[2]
\]
必有非平凡核。
\end{theorem}
\begin{proof}
由定理 \ref{thm:xi-window6-sm-envelope-center-2torsion-obstruction}，
\(
\dim_{\FF_2} Z(G_{\mathrm{SM}})[2]\le 2
\)。
而
\(
\mathcal C_{\mathrm{can}}\cong (\ZZ/2\ZZ)^3
\)
的 \(\FF_2\)-维数为 \(3\)。
故线性映射 \(f_*\) 的秩至多 \(2\)，其核维数至少 \(1\)，即核非平凡。证毕。
\end{proof}

\begin{proposition}[sheet parity 在 \(m=6\) 的严格锚点性]\label{prop:xi-window6-sheet-parity-anchor-m6}
在 dyadic boundary uplift 口径下，取
\[
\Delta_m^{\mathrm{bdry}}
:=\{n-V_m(w):\ n\in(\Fold_m^{\mathrm{bin}})^{-1}(w)\}
\]
（对 boundary \(w\) 不依赖于具体词）。
核验给出
\[
\Delta_6^{\mathrm{bdry}}=\{0,34\},\qquad
\Delta_7^{\mathrm{bdry}}=\{0,55,89\},\qquad
\Delta_8^{\mathrm{bdry}}=\{0,89,144\}.
\]
在“parity 由 uplift 数据无需额外选择即可恢复”的准则下，该结构仅在 \(m=6\) 成立。
\end{proposition}
\begin{proof}
若 parity 可由 \(\Delta_m^{\mathrm{bdry}}\) 唯一恢复，则每个 boundary 纤维必须存在唯一二元配对规则；
这要求非零 uplift 增量唯一。
当 \(m=6\) 时，\(\Delta_6^{\mathrm{bdry}}\) 仅含一个非零元 \(34\)，且 boundary 预像层数为 \(2\)，配对唯一。
当 \(m=7,8\) 时，\(\Delta_m^{\mathrm{bdry}}\) 各含两个不同非零元且预像层数为 \(3\)；
任意二元化都需额外选择（例如忽略某一层），不再由数据唯一决定。证毕。
\end{proof}

\begin{theorem}[左--右交换下的规范中心模与不可下降通道]\label{thm:xi-window6-lr-swap-canonical-center-mod}
在 \(\mathcal C_{\mathrm{can}}\cong(\ZZ/2\ZZ)^3\) 中取基
\[
a:=(-I_4,1,1),\qquad b:=(1,-I_2,1),\qquad c:=(1,1,-I_2).
\]
令 \(\iota\) 为交换 \(SU(2)_L\leftrightarrow SU(2)_R\) 的自同构，则
\[
\iota(a)=a,\qquad \iota(b)=c,\qquad \iota(c)=b.
\]
定义满同态
\[
\pi_{\mathrm{LR}}:\mathcal C_{\mathrm{can}}\twoheadrightarrow \ZZ/2\ZZ,\qquad
\pi_{\mathrm{LR}}(\alpha a+\beta b+\gamma c):=\beta+\gamma.
\]
则
\[
\ker(\pi_{\mathrm{LR}})=\mathcal C_{\mathrm{can}}^{\iota}
\cong (\ZZ/2\ZZ)^2.
\]
因而任何要求“左--右交换在中心层不可见”的下降机制，必然至少压缩一个规范 \(\ZZ_2\) 模通道。
\end{theorem}
\begin{proof}
\(\pi_{\mathrm{LR}}\) 显然为群同态且满。
\[
\pi_{\mathrm{LR}}(\alpha,\beta,\gamma)=0
\iff \beta=\gamma.
\]
另一方面，
\[
\iota(\alpha,\beta,\gamma)=(\alpha,\gamma,\beta),
\]
故
\(
(\alpha,\beta,\gamma)\in\mathcal C_{\mathrm{can}}^{\iota}
\iff \beta=\gamma
\)。
两条件等价，故核恰为不动子群，且其维数为 \(2\)。证毕。
\end{proof}

\begin{conjecture}[dyadic boundary uplift 层数在 \(m=6\) 的唯一退化]\label{conj:xi-window6-bdry-uplift-unique-degeneration}
设 \(\Delta_m^{\mathrm{bdry}}\) 为 boundary uplift 的增量集合。猜想对所有 \(m\ge 7\) 有
\[
\Delta_m^{\mathrm{bdry}}=\{0,F_{m+3},F_{m+4}\},
\]
从而 boundary 覆盖层数恒为 \(3\)。
而 \(m=6\) 为唯一退化分辨率：
\[
\Delta_6^{\mathrm{bdry}}=\{0,F_9\}=\{0,34\},
\]
即唯一仅含单一非零增量的二层情形。
\end{conjecture}

\begin{remark}[可复现核验脚本]\label{rem:xi-window6-center-sheetflip-ps-sm-audit-script}
本段代数与组合核验由脚本
\texttt{scripts/exp\_xi\_window6\_center\_sheetflip\_ps\_sm\_audit.py}
给出，证书文件为
\texttt{artifacts/export/xi\_window6\_center\_sheetflip\_ps\_sm\_audit.json}。
\end{remark}

